% Chapter 18: Frontiers and Open Problems
% Theme: WHAT WE DO NOT KNOW: question-driven structure.
% Each section leads with an open problem, shows what is known,
% identifies the gap, and asks what would close it.
% Concepts: compression_scheme, margin_theory, kolmogorov_complexity,
% algorithmic_probability, universal_learning, computational_hardness,
% differential_privacy_learning, active_learning, transfer/meta/continual.

\chapter{Frontiers and Open Problems}
\label{ch:frontiers}

Every preceding chapter has presented established results: definitions,
characterization theorems, separation witnesses, and tight bounds.  This
chapter presents none of these.  Its subject is the boundary of what we know,
and its organizing principle is the open question rather than the proven
theorem.

Each section leads with a question to which no complete answer is known,
presents the best partial results, identifies the precise gap between what is
known and what is conjectured, and asks what kind of argument might close it.

The problems chosen are those that
connect most tightly to the concept graph of this book: the compression
conjecture (Chapter~\ref{ch:compression}), deep network generalization
(Chapter~\ref{ch:generalization}), universal learning beyond countable classes
(Chapter~\ref{ch:universal}), computational--statistical tradeoffs
(Chapter~\ref{ch:computational}), and the role of Kolmogorov complexity in
learning.  Other important open problems (online multiclass learnability,
private agnostic learning, quantum sample complexity gaps) are noted but not
developed.

\medskip

The chapter has five sections.

\begin{enumerate}[leftmargin=*,itemsep=2pt]
\item \textbf{The Compression Conjecture} (\Cref{sec:compression-conjecture}):
  does $\VC(\mathcal{H}) = d$ imply a compression scheme of size $O(d)$?
\item \textbf{Deep Learning and Generalization}
  (\Cref{sec:deep-generalization}): why do overparameterized networks
  generalize?
\item \textbf{Universal Learning Beyond Countable Classes}
  (\Cref{sec:universal-frontiers}): does the trichotomy extend?
\item \textbf{Computational--Statistical Tradeoffs}
  (\Cref{sec:comp-stat-tradeoffs}): when does computational hardness prevent
  statistically optimal learning?
\item \textbf{Kolmogorov Complexity and the Ideal Learner}
  (\Cref{sec:kolmogorov-learning}): can the uncomputable yield computable
  insight?
\end{enumerate}


% ================================================================
% SECTION 1: THE COMPRESSION CONJECTURE
% ================================================================

\section{The Compression Conjecture}
\label{sec:compression-conjecture}

\begin{openprob}[Littlestone--Warmuth Compression Conjecture]
\label{op:compression-frontiers}
Let $\mathcal{H} \subseteq \{0,1\}^X$ with $\VC(\mathcal{H}) = d < \infty$.
Does there exist a sample compression scheme for $\mathcal{H}$ of size
$O(d)$?
\end{openprob}

This conjecture, stated by Littlestone and Warmuth in 1986, is among the
oldest major open problems in computational learning theory.  It asks whether
the sample complexity of PAC learning (which is $\Theta(d/\varepsilon)$) is
explained by compression: whether every learnable class can be learned by
storing only $O(d)$ examples from the training set and reconstructing a
consistent hypothesis from those examples alone.

\subsection{What is Known}

The positive direction is settled in a weak sense.

\begin{itemize}[leftmargin=*]
\item \textbf{Exponential compression exists.}
  Moran and Yehudayoff~\cite{MoranYehudayoff2016} proved that every class
  with $\VC(\mathcal{H}) = d$ admits a compression scheme of size
  $2^{O(d)}$.  The proof uses the Radon partition structure of finite VC
  classes, but the bound is exponentially larger than the conjecture predicts.

\item \textbf{Special cases are resolved.}  Maximum classes (those achieving
  the Sauer--Shelah bound $\growth_\mathcal{H}(n) = \binom{n}{\leq d}$)
  admit compression of size exactly $d$, as shown by Floyd and
  Warmuth~\cite{FloydWarmuth1995}.  Intersection-closed classes admit
  compression of size $d$~\cite{Helmbold2018}, and classes of small recursive
  teaching dimension compress to that dimension.

\item \textbf{Compression implies learning.}  The converse direction is
  established: a compression scheme of size $k$ yields PAC learning with
  sample complexity $O\bigl(\frac{k}{\varepsilon}\log\frac{k}{\varepsilon}
  + \frac{1}{\varepsilon}\log\frac{1}{\delta}\bigr)$.  Thus compression
  of size $O(d)$ would give an alternative proof of the fundamental theorem
  of PAC learning, with the added structural insight that the hypothesis
  depends on only $O(d)$ training points.
\end{itemize}

\begin{historical}
The compression conjecture predates the fundamental theorem of statistical
learning.  Littlestone and Warmuth formulated it in 1986, three years before
Blumer, Ehrenfeucht, Haussler, and Warmuth~\cite{BlumerEhrenfeuchtHausslerWarmuth1989}
completed the characterization of PAC learnability through VC dimension.
The conjecture was motivated by the observation that specific learnable
classes (halfspaces, rectangles, decision lists) all admit small compression
schemes, and the belief that this regularity was a property of finite-VC
classes as such.  Four decades later, the general case remains open.
\end{historical}

\subsection{Where the Gap Is}

The gap is between $O(d)$ and $2^{O(d)}$.  No class is known to require
compression size $\omega(d)$, but no proof technique is known to achieve
$o(2^d)$ in general.

The difficulty is structural.  The Moran--Yehudayoff proof proceeds through
a topological argument (the staircase lemma) that inherently produces
exponential-size compressions.  Improving the bound requires a fundamentally
different approach, one that exploits the combinatorial structure of VC
classes more directly.

Several partial strategies have been explored:

\begin{itemize}[leftmargin=*]
\item \textbf{Unlabeled compression.}  Allowing the compression to store
  unlabeled points (without their labels) does not help: the reconstruction
  map must still determine labels, and known reductions show that unlabeled
  compression of size $k$ implies labeled compression of size $O(k)$.

\item \textbf{Algebraic approaches.}  For classes definable by polynomial
  threshold functions, algebraic arguments yield compression of size $O(d)$.
  But this exploits algebraic structure absent from general VC classes.

\item \textbf{Recursive approaches.}  Decomposing a class into subclasses
  of smaller VC dimension and compressing each recursively is a natural
  strategy, but known decompositions (e.g., via Helly-type theorems) produce
  too many subclasses, losing the linear bound.
\end{itemize}

\subsection{What Would Close the Gap}

A resolution likely requires one of the following:

\begin{enumerate}[leftmargin=*]
\item A new structural theorem about VC classes that controls the geometry
  of their restrictions to finite sets more tightly than Sauer--Shelah.
\item A connection to a different branch of combinatorics (e.g., matroid
  theory or extremal set theory) that provides the right decomposition lemma.
\item A counterexample: a family of classes with VC dimension $d$ and
  no compression scheme of size $o(d^{1+c})$ for some $c > 0$.  Even a
  superlinear lower bound would be a breakthrough.
\end{enumerate}

\begin{remark}
The compression conjecture illustrates a recurring theme in this chapter:
the gap between \emph{existential} and \emph{constructive} knowledge.  We
know that every finite-VC-dimension class is PAC learnable (Chapter~\ref{ch:pac}).
We know that compression of size $2^{O(d)}$ exists
(Chapter~\ref{ch:compression}).  What we do not know is whether these two
facts are reflections of the same underlying structure at the right
quantitative scale.
\end{remark}


% ================================================================
% SECTION 2: DEEP LEARNING AND GENERALIZATION
% ================================================================

\section{Deep Learning and Generalization}
\label{sec:deep-generalization}

\begin{openprob}[Generalization in Overparameterized Networks]
\label{op:deep-gen}
Let $\mathcal{H}$ be the class of functions computed by a neural network
with $p \gg n$ parameters, trained by gradient descent on $n$ samples.
Why does the trained network generalize, i.e., why is $\risk(\hat{h})
\approx \emprisk(\hat{h})$, despite $\VC(\mathcal{H}) \geq p$?
\end{openprob}

This is among the most practically urgent open problems in learning theory.
Classical generalization bounds predict that a model with more parameters than
training examples should overfit catastrophically.  Deep networks routinely
violate this prediction.

\subsection{What is Known}

Several partial explanations exist, none fully satisfactory.

\begin{itemize}[leftmargin=*]
\item \textbf{Margin-based bounds.}  Bartlett, Foster, and
  Telgarsky~\cite{BartlettFosterTelgarsky2017} proved that for a
  depth-$L$ network with weight matrices $W_1, \ldots, W_L$, the
  generalization error is controlled by
  \[
    \frac{1}{n}\biggl(\frac{\prod_{i=1}^L \|W_i\|_\sigma}{\gamma}\biggr)^2
    \cdot \sum_{i=1}^L \frac{\|W_i\|_F^2}{\|W_i\|_\sigma^2}
  \]
  where $\gamma$ is the margin, $\|\cdot\|_\sigma$ is the spectral norm,
  and $\|\cdot\|_F$ is the Frobenius norm.  This bound is independent of the
  number of parameters $p$ and depends instead on the \emph{spectral
  complexity} of the learned weights.

\item \textbf{PAC-Bayes bounds.}  Treating the trained weights as a
  posterior distribution (e.g., by adding Gaussian noise), PAC-Bayes
  bounds (\Cref{ch:extensions}) can yield non-vacuous generalization
  certificates for specific trained networks.  Dziugaite and
  Roy~\cite{DziugaiteRoy2017} demonstrated this empirically.

\item \textbf{Algorithmic stability.}  SGD with bounded iterations satisfies
  uniform stability with $\beta = O(1/n)$ under convexity assumptions.
  For non-convex losses, Hardt, Recht, and Singer~\cite{HardtRechtSinger2016}
  showed that early stopping provides stability guarantees, but the bounds
  depend on the number of SGD steps and become vacuous for the training
  durations used in practice.

\item \textbf{Compression-based arguments.}  If the effective hypothesis
  after training can be compressed to $k \ll p$ bits, then compression
  bounds (Chapter~\ref{ch:compression}) apply.  Lottery ticket
  observations (that trained networks contain sparse subnetworks matching
  full performance) provide indirect evidence for compression, but no
  rigorous connection to sample complexity has been established.

\item \textbf{Machine-checked certificates.}  For the specific function class
  of a multi-head attention block, the generalization certificate of
  \Cref{ch:generalization} has been formalized: the executed
  (finite-precision) network's true risk exceeds its empirical risk by more
  than an explicit Dudley-capacity budget only on a sample event of
  probability at most~$\delta$.\formal{TLT.CertifiedGeneralization}  The budget
  is discharged from the layer norms and depth of the
  stack.\formal{TLT.config_capacity_law}  Like its hand-proved cousins, it
  grows with depth and norm and is not yet non-vacuous on large models.
\end{itemize}

\begin{computational}
\textbf{The double descent curve.}  Consider interpolation threshold: the
point where the number of parameters $p$ equals the number of training
examples $n$.  Classical theory predicts that test error is U-shaped
as a function of model complexity, with minimum at $p \approx n$.
Empirically, the test error \emph{decreases again} for $p \gg n$,
forming a double descent curve.  This was systematically documented
by Belkin et al.~\cite{BelkinHsuMaMandal2019}.

The double descent phenomenon is observed across several model families
(linear models, random features, deep networks)~\cite{BelkinHsuMaMandal2019}.
The classical complexity-based bounds (VC, Rademacher, and the parameter-count
bounds of \Cref{ch:generalization}) are monotone in model complexity and so
cannot predict the second descent; whether \emph{some} data-dependent bound can
is itself part of \Cref{op:deep-gen}.  The relevant notion of complexity for
overparameterized models is not a simple function of parameter count.
\end{computational}

\subsection{Where the Gap Is}

The gap has three aspects.

\begin{enumerate}[leftmargin=*]
\item \textbf{Quantitative vacuity.}  Most norm- and parameter-count-based
  bounds, when evaluated on practical networks, predict generalization error
  close to~$1$; that is, they are vacuous.  The PAC-Bayes certificates of
  Dziugaite and Roy~\cite{DziugaiteRoy2017} are a notable exception, giving
  non-vacuous numbers for specific small networks, but no bound is yet both
  non-vacuous and tight on large benchmark models.

\item \textbf{Architecture dependence.}  Existing bounds treat networks as
  generic function classes parameterized by weight norms.  They do not
  explain why certain architectures (e.g., ResNets, Transformers)
  generalize better than others with the same parameter count and weight
  norms.

\item \textbf{Optimization dependence.}  The generalization behavior of deep
  networks depends heavily on the optimization algorithm (SGD vs.\ Adam),
  learning rate schedule, batch size, and initialization.  No existing
  framework captures how the optimization trajectory selects hypotheses
  with low population risk from the vast set of interpolating solutions.
\end{enumerate}

\subsection{What Would Close the Gap}

The field needs a complexity measure that simultaneously accounts for the
architecture, the full optimization trajectory, and the data distribution.

Possible directions include:
\begin{itemize}[leftmargin=*]
\item \textbf{Implicit regularization theory:} proving that gradient descent
  on specific architectures converges to solutions with low complexity
  in a suitable measure (e.g., low rank, low spectral norm product, or
  low description length).
\item \textbf{Distribution-dependent bounds:} replacing uniform convergence
  (which bounds $\sup_D |\risk - \emprisk|$) with bounds that depend on the
  specific data distribution, potentially through conditional mutual
  information (Steinke--Zakynthinou, \Cref{ch:generalization}).
\item \textbf{Feature learning theory:} understanding how networks transform
  input representations during training, and why the learned features
  support generalization.
\end{itemize}


% ================================================================
% SECTION 3: UNIVERSAL LEARNING BEYOND COUNTABLE CLASSES
% ================================================================

\section{Universal Learning Beyond Countable Classes}
\label{sec:universal-frontiers}

\begin{openprob}[Universal Learning for Uncountable Classes]
\label{op:universal-uncountable}
The trichotomy theorem of Bousquet et
al.~\cite{BousquetHannekeMoranVanHandelYehudayoff2021} characterizes
universal learning rates for classes $\mathcal{H}$ with $|\mathcal{H}| \geq 3$.
Does an analogous characterization hold for uncountable classes of functions
$f \colon X \to [0,1]$ under general loss functions?
\end{openprob}

\subsection{What is Known}

The trichotomy theorem (\Cref{ch:universal}) establishes that for hypothesis
classes of size at least 3, the optimal universal learning rate is exactly one
of three types:
\begin{enumerate}[leftmargin=*,nosep]
\item Exponential ($e^{-n}$): if and only if the Littlestone dimension is
  finite.
\item Linear ($1/n$): if and only if the Littlestone dimension is infinite
  but no infinite VCL tree exists.
\item Arbitrarily slow: if and only if an infinite VCL tree exists.
\end{enumerate}
This characterization applies to binary classification with the $0$-$1$ loss.

For regression and general losses, partial results exist:
\begin{itemize}[leftmargin=*]
\item Hanneke, Moran, and others have extended the trichotomy to multiclass
  settings with finite label spaces, showing that the DS dimension
  (\Cref{ch:extensions}) plays a role analogous to the Littlestone dimension.
\item For real-valued regression under squared loss, the fat-shattering
  dimension at all scales characterizes PAC learnability, but the universal
  learning rate structure is not yet characterized.
\end{itemize}

\subsection{Where the Gap Is}

The gap has two components:

\begin{enumerate}[leftmargin=*]
\item \textbf{The loss function.}  The trichotomy proof relies on the
  discreteness of $0$-$1$ loss.  For continuous losses (squared loss,
  absolute loss, logistic loss), the proof technique does not directly
  apply.  The VCL tree, which is the combinatorial object governing the
  third regime, has no established analogue for continuous losses.

\item \textbf{The function class.}  For uncountable classes (e.g., all
  Lipschitz functions, all bounded-variation functions), measurability
  issues arise: the supremum of uncountably many random variables may not
  be measurable without additional regularity assumptions.  The universal
  learning framework requires that the learning rate hold for
  \emph{all} realizable distributions, and handling uncountable classes
  demands careful measure-theoretic foundations.
\end{enumerate}

\subsection{What Would Close the Gap}

Two ingredients appear necessary:
\begin{itemize}[leftmargin=*]
\item A combinatorial object for continuous losses that plays the role of
  the Littlestone tree and VCL tree in the binary case, a ``scale-sensitive
  game tree'' that captures the sequential complexity of the class at all
  accuracy levels.
\item A measurability framework (e.g., based on metric entropy or covering
  numbers) that allows the universal learning definition to extend cleanly
  to uncountable classes without requiring separability or other ad hoc
  assumptions.
\end{itemize}


% ================================================================
% SECTION 4: COMPUTATIONAL--STATISTICAL TRADEOFFS
% ================================================================

\section{Computational--Statistical Tradeoffs}
\label{sec:comp-stat-tradeoffs}

\begin{openprob}[Computational--Statistical Gap]
\label{op:comp-stat}
For which concept classes $\mathcal{C}$ does a gap exist between the
information-theoretic sample complexity of PAC learning $\mathcal{C}$ and
the sample complexity achievable by any polynomial-time algorithm?
Can such gaps be characterized combinatorially?
\end{openprob}

The results of Chapter~\ref{ch:computational} show that computational
hardness can prevent efficient PAC learning even when the VC dimension is
finite.  But those results are conditional: they assume the hardness of
specific cryptographic problems (factoring, LWE, discrete cube root).  The
deeper question is whether the computational--statistical gap is a structural
feature of learning or an artifact of our inability to prove $P \neq NP$.

\subsection{What is Known}

\begin{itemize}[leftmargin=*]
\item \textbf{Cryptographic hardness.}  Kearns and
  Valiant~\cite{KearnsValiant1994} showed that under the Discrete Cube Root
  Assumption, polynomial-size circuits are not PAC learnable in polynomial
  time, despite having polynomial VC dimension.  The gap is between
  $O(d/\varepsilon)$ samples (information-theoretically sufficient) and
  ``no polynomial-time algorithm suffices'' (computationally).

\item \textbf{Statistical query lower bounds.}  The SQ dimension provides
  \emph{unconditional} computational lower bounds within the SQ model
  (\Cref{ch:computational}).  Parities over $\{0,1\}^n$ have
  $\VC = n$ but $\SQdim = 2^n$: any SQ algorithm needs exponentially
  many queries.  Yet Gaussian elimination learns parities in polynomial
  time; it is simply not an SQ algorithm.  This shows that the SQ model
  captures one type of computational constraint but not all.

\item \textbf{Planted clique and related problems.}  A growing body of work
  studies problems where the information-theoretic threshold differs from the
  conjectured computational threshold: planted clique, sparse PCA, community
  detection.  For planted clique, a planted clique of size $\Theta(\sqrt{n})$
  is the conjectured threshold for polynomial-time detection, while one of
  size $\sim 2\log_2 n$ is already detectable information-theoretically.  None
  of these gaps is proven unconditionally; the evidence comes from restricted
  models (sum-of-squares, low-degree polynomials) rather than from
  unconditional separations.

\item \textbf{The barrier of Applebaum--Barak--Xiao.}  Applebaum, Barak, and
  Xiao~\cite{ApplebaumBarakXiao2008} showed that one cannot prove
  ``$\mathrm{P} \neq \mathrm{NP}$ implies hard PAC learning'' via standard
  (black-box) reductions.  Cryptographic or average-case assumptions are
  inherently necessary for hardness results in PAC learning.  This is a
  meta-theoretic constraint on the kind of hardness proofs that are possible.
\end{itemize}

\begin{historical}
The computational--statistical tradeoff emerged as a distinct research
direction in the 2010s, though its roots go back to Kearns and Valiant's
1994 cryptographic hardness results.  The key conceptual shift was the
recognition that sample complexity and computational complexity are not
independent axes of difficulty, but interact in subtle ways.  A problem
may have low sample complexity (few examples suffice) but high
computational complexity (no efficient algorithm can use those examples),
or vice versa.  The SQ model, introduced by Kearns~\cite{Kearns1998},
provided the first framework where computational lower bounds could be
proved unconditionally, but only within a restricted model of computation.
\end{historical}

\subsection{Where the Gap Is}

The fundamental gap is proof-theoretic: we lack techniques to prove
unconditional computational lower bounds for PAC learning.

\begin{enumerate}[leftmargin=*]
\item \textbf{No unconditional separations.}  Every known example of a
  computational--statistical gap relies on a complexity-theoretic assumption.
  Proving an unconditional gap would require proving $P \neq NP$ or something
  similarly difficult.

\item \textbf{No combinatorial characterization.}  For information-theoretic
  learnability, we have a complete combinatorial characterization: finite VC
  dimension.  For computationally efficient learnability, no analogous
  characterization exists.  We do not know what combinatorial property of a
  concept class determines whether it is efficiently learnable.

\item \textbf{Model dependence.}  Lower bounds proved in the SQ model, the
  low-degree polynomial model, or the sum-of-squares hierarchy apply only to
  those specific computational models.  An algorithm outside the model
  (like Gaussian elimination for parities) may circumvent the lower bound.
  No model of computation is known to be simultaneously natural, powerful
  enough to capture practical algorithms, and amenable to unconditional lower
  bounds.
\end{enumerate}

\subsection{What Would Close the Gap}

\begin{itemize}[leftmargin=*]
\item \textbf{A natural proof barrier for learning.}  Razborov and
  Rudich's natural proofs barrier explains why circuit lower bounds are hard.
  An analogous barrier for learning (explaining why computational lower
  bounds for PAC learning are hard) would clarify the structural limits of
  current proof techniques.
\item \textbf{A unifying computational model.}  A model of bounded
  computation that captures SQ algorithms, low-degree methods,
  sum-of-squares, and gradient descent, while admitting provable lower
  bounds, would provide a natural home for computational--statistical
  tradeoff results.
\item \textbf{Conditional characterizations.}  Even under assumptions
  like $P \neq NP$ or the hardness of LWE, a characterization of
  efficiently learnable classes (analogous to the fundamental theorem for
  information-theoretic learnability) would be a major advance.
\end{itemize}


% ================================================================
% SECTION 5: KOLMOGOROV COMPLEXITY AND THE IDEAL LEARNER
% ================================================================

\section{Kolmogorov Complexity and the Ideal Learner}
\label{sec:kolmogorov-learning}

\begin{openprob}[Computable Approximations to Solomonoff Induction]
\label{op:solomonoff-approx}
Solomonoff's universal prior $M(x) = \sum_{p:\, U(p)=x} 2^{-|p|}$ defines
an ideal but uncomputable learner.  To what extent can computable
approximations to $M$ achieve the learning-theoretic properties of the ideal
learner?
\end{openprob}

This section concerns the deepest stratum of the theory: the connection
between algorithmic information theory and learning.

\subsection{What is Known}

\begin{itemize}[leftmargin=*]
\item \textbf{Solomonoff's completeness theorem.}  For any computable
  measure $\mu$ on infinite binary sequences, the universal semimeasure $M$
  satisfies
  \[
    \sum_{n=1}^\infty \E_{x_{1:n} \sim \mu}
    \bigl[D_{\mathrm{KL}}(\mu(\cdot \mid x_{1:n}) \,\|\,
    M(\cdot \mid x_{1:n}))\bigr] < \infty.
  \]
  That is, the cumulative prediction error of $M$ relative to any
  computable environment is bounded.  This makes $M$ the ``ideal learner'':
  it converges to the true distribution at a rate dominated by the
  Kolmogorov complexity of the environment.

\item \textbf{MDL and MML as approximations.}  The minimum description
  length principle (Rissanen~\cite{Rissanen1984}) and minimum message length
  (Wallace and Boulton~\cite{WallaceBoulton1968}) are computable
  approximations to the Kolmogorov-optimal encoding.  Both embody a version
  of Occam's razor: prefer the hypothesis with the shortest total description
  (model + data given model).

\item \textbf{Occam bounds via Kolmogorov complexity.}  Li, Tromp, and
  Vit\'anyi~\cite{LiVitanyi2008} showed that replacing description length
  with Kolmogorov complexity in Occam-style PAC bounds yields tighter
  (but uncomputable) generalization guarantees.  The bound states that
  the generalization error of a hypothesis $h$ is controlled by
  $K(h)/n$, where $K(h)$ is the Kolmogorov complexity of $h$ and $n$
  is the sample size.

\item \textbf{Algorithmic probability and universal learning.}  Universal
  learning allows distribution-dependent rates; Solomonoff's prior achieves
  distribution-dependent convergence.  But the universal learning framework is
  information-theoretic (it asks only about sample complexity), while
  Solomonoff's prior is algorithmic (it asks about description complexity).
  The precise relationship between these two notions of universality is open.
\end{itemize}

\begin{historical}
Solomonoff's 1964 paper~\cite{Solomonoff1964} predates Gold's 1967
formalization of identification in the limit and Valiant's 1984
introduction of PAC learning.  It is among the first formal proposals
for a universal learning algorithm, though it was formulated in the language
of algorithmic information theory rather than learning theory.  The
intellectual genealogy runs:
Solomonoff~(1964) $\to$ Kolmogorov complexity~(1965) $\to$ MDL~(1978)
$\to$ Occam bounds in PAC~(1987) $\to$ compression schemes~(1986).
The compression conjecture of \Cref{sec:compression-conjecture} can be
viewed as asking whether the Occam principle (short description implies
good generalization) holds at the tightest possible quantitative level.
\end{historical}

\subsection{Where the Gap Is}

\begin{enumerate}[leftmargin=*]
\item \textbf{Uncomputability.}  Kolmogorov complexity and Solomonoff's
  prior are uncomputable.  Every computable approximation introduces a
  gap: MDL uses a fixed model class rather than all programs; bounded
  Solomonoff approximations (running programs for at most $t$ steps)
  introduce a time--complexity tradeoff that is poorly understood.

\item \textbf{No convergence rate for approximations.}  For the ideal
  learner $M$, the convergence rate is known: the expected KL divergence
  is $O(K(\mu)/n)$ where $K(\mu)$ is the complexity of the true
  environment.  For computable approximations, no analogous rate is
  established that holds uniformly over all computable environments.

\item \textbf{The gap between prediction and classification.}  Solomonoff's
  framework concerns \emph{sequence prediction}: predicting the next bit.
  PAC learning concerns \emph{concept identification}: finding a hypothesis
  close to the target.  The formal relationship between these two tasks
  is well understood in specific settings (e.g., the realizable case with
  i.i.d.\ data), but the general connection (especially in the agnostic
  or adversarial setting) remains underdeveloped.
\end{enumerate}

\subsection{What Would Close the Gap}

\begin{itemize}[leftmargin=*]
\item \textbf{Time-bounded Solomonoff induction.}  A theory of
  $M^t(x) = \sum_{p:\, U(p)=x,\; \mathrm{time}(p) \leq t} 2^{-|p|}$
  with provable convergence rates as a function of both $n$ (sample size)
  and $t$ (computation time) would bridge the gap between the ideal and
  the computable.  Such a theory would need to handle the fact that
  $M^t$ is not a semimeasure for finite $t$, and that the dominant programs
  may have very long running times.

\item \textbf{Kolmogorov complexity as a complexity measure for PAC
  learning.}  Formalizing $K(\mathcal{H})$ (the description complexity of
  a hypothesis class) as a complexity measure in the sense of
  Chapter~\ref{ch:dimensions} (with upper and lower bound relationships
  to VC dimension, Rademacher complexity, and compression size) would
  integrate algorithmic information theory into the concept graph of this
  book.

\item \textbf{A computable universal learner.}  An algorithm that, for
  every computable concept class $\mathcal{C}$ and every realizable
  distribution $D$, achieves PAC learning with sample complexity depending
  on $K(\mathcal{C})$, would be a computable analogue of the Solomonoff
  ideal.  Whether such an algorithm exists is itself an open question,
  closely related to the universal learning framework of
  Chapter~\ref{ch:universal}.
\end{itemize}


% ================================================================
% SECTION 6: FURTHER FRONTIERS
% ================================================================

\section{Further Frontiers}
\label{sec:further-frontiers}

The five problems above are those most tightly connected to the conceptual
framework of this book.  We note several other directions where the theory is
actively expanding, each involving mathematical structures that do not reduce
to the types established in Chapters~1--17.

\begin{openprob}[Private agnostic learning]
\label{op:private}
Differential privacy constrains the learner: the output must not reveal too
much about any individual example.  Kasiviswanathan et
al.~\cite{KasiviswanathanLeeNissimRaskhodnikovaSmith2011} showed that private
PAC learning under \emph{pure} differential privacy has sample complexity
$\Theta\bigl(d/\varepsilon + \log|\mathcal{X}|/(\varepsilon\alpha)\bigr)$, and
for \emph{approximate} privacy the Littlestone dimension characterizes private
learnability~\cite{AlonLivniMalliarisMoran2019,BunLivniMoran2020}, linking
online learning to privacy.  What is the tight sample complexity of private
\emph{agnostic} learning?
\end{openprob}

\begin{openprob}[Active learning under agnostic noise]
\label{op:active}
An active learner queries labels selectively, and the label complexity can be
exponentially smaller than passive sample complexity \emph{for favorable
classes and distributions} (for instance, halfspaces under a margin).
Hanneke~\cite{Hanneke2014} showed that the disagreement coefficient controls
label complexity for general classes.  What is the tight characterization of
active learning under \emph{agnostic} noise, and what is its computational
complexity?
\end{openprob}

\begin{openprob}[Transfer, meta-, and continual learning]
\label{op:meta}
Formal frameworks for learning across related tasks exist (Baxter's model of
learning to learn~\cite{Baxter2000}, and the PAC-Bayes approach to
meta-learning), but no characterization theorem analogous to the VC
characterization is known.  Can the combinatorial objects governing multi-task
learning be expressed through the single-task dimensions of this book, or are
genuinely new measures required?
\end{openprob}

\begin{openprob}[Online multiclass learnability]
\label{op:online-multiclass-frontiers}
The Littlestone dimension characterizes online binary learnability, but the
Littlestone dimension of the binary reductions does not always capture the
optimal multiclass mistake bound.  What combinatorial parameter characterizes
online multiclass learnability, and how does it relate to the DS dimension
(\Cref{ch:extensions})?
\end{openprob}

\begin{openprob}[Quantum sample complexity]
\label{op:quantum}
In quantum PAC learning the learner receives quantum superpositions of labeled
examples.  For binary classification the quantum and classical sample
complexities are equal up to constant factors, but for other tasks and for
\emph{time} complexity the relationship is unsettled.  Characterize, in terms
of classical complexity measures, exactly when quantum examples reduce the cost
of learning.
\end{openprob}

\begin{openprob}[A non-vacuous \emph{verified} bound for transformers]
\label{op:certified-transformer}
\Cref{sec:deep-generalization} notes that a machine-checked generalization
certificate already exists for an attention block
(\leanname{config_capacity_law}), but its Dudley budget grows with depth and
norm.  Can a generalization bound for a realistic transformer be \emph{both}
machine-checked and non-vacuous on a benchmark, thereby closing \Cref{op:deep-gen}?
\end{openprob}

\bigskip

\begin{computational}
\textbf{The shape of the frontier.}
The open problems of this chapter are not isolated.  They form a connected
subgraph in the concept graph, linked by shared dependencies.  The
compression conjecture involves the same VC dimension that governs PAC
learning, which connects to the computational hardness that creates
computational--statistical gaps, which connects to the SQ dimension, which
connects to the Kolmogorov complexity that underlies the ideal learner.
The deep learning generalization problem is linked to compression (through
lottery tickets), to margin theory (through spectral bounds), and to
algorithmic stability (through SGD analysis).  Universal learning connects
to the Littlestone dimension and therefore to private learning.

This interconnection is not coincidental.  A resolution of the compression
conjecture would yield insight into generalization bounds; a combinatorial
characterization of efficient learnability would require new structural
theorems about VC classes; a computable approximation to Solomonoff induction
would connect to both compression and universal learning.  Resolving any one
of these problems would reshape the others.
\end{computational}

\bigskip

This chapter has surveyed what we do not know.  The theory developed in
Chapters~1--17 provides a rigorous, interconnected framework for what we do
know: characterization theorems, tight bounds, separation results, and
structural analogies.  The open problems of this chapter mark the places
where that framework reaches its current limits, and where the next results
will be found.
